\scanpage{26}
1. $X$ espace top.

$Y$ sous-espace ``non singulier'' i.e. variété

$X$ est dit ``\uline{équisingulier le long de} $Y$'' si $\forall y\in Y$, \deleted{\unsure{alors}} $\exists$ espace pointé $(N,a)$, \deleted{et} un voisinage ouvert $U$ de $y$ dans $X$, et un homéom
\diagram{\begin{tikzcd}[column sep=1.4em]
{} & U \arrow[r,phantom,"\simeq"] & N\times\unsure{Y|V}_U & (Y_U=Y\cap U) \\
\text{avec commutativité:} & Y_U \arrow[u,hook] \arrow[r,"\simeq","\mathrm{can}"'] & a\times Y_U \arrow[u,hook] & {}
\end{tikzcd}}

On dit que $x\in X$ est une \uline{singularité élémentaire} \inserted{(ou \deleted{\unsure{spéciale}} \inserted{immobile})} de $X$ si $\forall Y\subset X$, $Y$ sous-espace \deleted{\unsure{de $X$}} non-singulier de $X$, $Y\ni x$, $X$ équisingulier \deleted{\unsure{sur $Y$}} le long de $\unsure{Y}$, on a $\dim_xY=0$.

On dit que $x\in X$ est une \uline{singularité} \uline{$\nu$-mobil\unsure{e}\deleted{\unsure{ité}}} \deleted{\unsure{élémentaire}}, si $\exists Y\subset X$ sous-espace, $Y$ non sing., $Y\ni x$, $X$ \inserted{équi}singulier le long de $Y$, $\dim_xY\geq\nu$.

(N.B. on peut alors choisir $Y$ avec $\dim Y=\nu$).

\deleted{Ainsi, les singularités \unsure{élémentaires}}

On dit que $x$ a un \uline{ordre de mobilité} $\nu$ si $x$ est $\nu$-mobile mais non $(\nu+1)$-mobile. (Ainsi les singularités élémentaires \inserted{ou ``immobiles''} sont celles qui sont d'ordre de mobilité 0). \deleted{Donc les points} \deleted{\unsure{L'ens des points qui ont un ordre de} $\nu$-mobiles,} \inserted{Les points qui ont un ordre de mobilité} $\nu\geq\nu_0$ ($\nu_0$ donné) sont les p\textsuperscript{ts} $\nu_0$-mobiles qui ont un ordre de mobilité. (Mais si $X$ \deleted{\unsure{a}} est loc de dim finie, tt p\textsuperscript{t} \inserted{$x$} de $X$ a nécessairement un ordre de mobilité \inserted{$\nu(x)$}, qui est le sup des dim en $x$ des ss-espaces de $X$ \deleted{\unsure{singuliers}} non singuliers \deleted{qui sont} le long desquels $X$ est équisingulier)
\mspage{27}
